\documentclass[12pt]{article}

\usepackage[utf8]{inputenc}
\usepackage[T1]{fontenc}
\usepackage[brazil]{babel}

% Template oficial da SBC (sbc-template.sty e sbc.bst, do modelo "SBC Conferences"
% do Overleaf). Sem esses arquivos, o preâmbulo usa uma configuração equivalente
% para o documento compilar em qualquer instalação do LaTeX.
\newif\ifsbcoficial
\IfFileExists{sbc-template.sty}{\sbcoficialtrue}{\sbcoficialfalse}
\ifsbcoficial
  \usepackage{sbc-template}
\else
  \usepackage[a4paper,top=3.5cm,bottom=2.5cm,left=3cm,right=3cm]{geometry}
  \newcommand{\inst}[1]{\textsuperscript{#1}}
  \newcommand{\address}[1]{\gdef\enderecoautores{#1}}
  \makeatletter
  \renewcommand{\maketitle}{%
    \begin{center}
      {\Large\bfseries \@title\par}\vspace{1em}
      {\bfseries \@author\par}\vspace{0.6em}
      {\small \enderecoautores\par}
    \end{center}\vspace{1em}}
  \makeatother
  \renewenvironment{abstract}{\begin{quote}\small\textbf{\textit{Abstract.}}~\itshape}{\end{quote}}
  \newenvironment{resumo}{\begin{quote}\small\textbf{\textit{Resumo.}}~\itshape}{\end{quote}}
\fi
% "Autor [ano]" no texto. O template da SBC já define; senão, usa o \cite comum.
\providecommand{\citeasnoun}[1]{\cite{#1}}

\usepackage{times}
\usepackage{graphicx}
\usepackage{url}
\usepackage{booktabs}
\usepackage{array}
\usepackage{tikz}
\usetikzlibrary{positioning,arrows.meta,fit}

\sloppy

\title{GraphRAG e KAG sobre a história da computação:\\
perguntas multi-hop com um grafo de conhecimento no Neo4j}

\author{Victor Gabriel Basilio e Silva\inst{1}, Nicolas Brito Fernandes da Costa\inst{1}, Patrick Duarte Sardenberg\inst{1},\\
Gabriel Xavier Soares\inst{1}, Pedro Jorge Fernandes Medeiros de Oliveira\inst{1}}

\address{\inst{1}Centro Universitário de João Pessoa (UNIPÊ) -- João Pessoa, PB -- Brasil\\
Tendências em Ciência da Computação -- Prof. Ricardo Roberto de Lima -- 2026}

\begin{document}

\maketitle

\begin{abstract}
Large language models answer fluently but may invent facts, and retrieval-augmented
generation based only on vector search tends to fail on questions that require
combining facts from different documents. This paper presents a GraphRAG/KAG question
answering system about the history of computing, built from Portuguese Wikipedia
articles. Texts are split into passages, entities and relations are extracted by an
LLM under a predefined domain schema, and everything is stored as a labeled property
graph in Neo4j, where each relation points to the passage that supports it. Retrieval
combines vector search, multi-hop expansion over the graph and LLM-generated Cypher
queries, and the answer cites the paths and passages it used, with a confidence level
computed from retrieval signals. In an evaluation with 15 questions, GraphRAG correctly
answered 2 of 7 multi-hop questions, against 1 for vector RAG, and neither mode
produced a hallucinated answer.
\end{abstract}

\begin{resumo}
Modelos de linguagem respondem com fluência, mas podem inventar fatos, e a geração
aumentada por recuperação baseada só em busca vetorial tende a falhar em perguntas que
exigem combinar fatos de documentos diferentes. Este artigo apresenta um sistema
GraphRAG/KAG de perguntas e respostas sobre a história da computação, construído a
partir de artigos da Wikipédia em português. Os textos são divididos em trechos, as
entidades e relações são extraídas por um LLM sob um schema de domínio pré-definido, e
tudo é gravado como um grafo de propriedades rotuladas no Neo4j, em que cada relação
aponta para o trecho que a sustenta. A recuperação combina busca vetorial, expansão
multi-hop pelo grafo e consultas Cypher geradas por LLM, e a resposta cita os caminhos e
trechos usados, com um nível de confiança calculado a partir de sinais da recuperação.
Numa avaliação com 15 perguntas, o GraphRAG respondeu corretamente 2 de 7 perguntas
multi-hop, contra 1 do RAG vetorial, e nenhum dos modos produziu resposta alucinada.
\end{resumo}

% =====================================================================
\section{Introdução}

Modelos de linguagem de grande porte (LLMs) produzem texto fluente, mas podem afirmar
fatos que não constam de nenhuma fonte, fenômeno conhecido como alucinação
\cite{ji2023survey}. A geração aumentada por recuperação (RAG) reduz esse risco ao
entregar ao modelo, junto com a pergunta, trechos de documentos recuperados de uma base
de conhecimento, e ao pedir que a resposta se apoie neles \cite{lewis2020rag}.

Na forma mais comum, a recuperação é uma busca vetorial: a pergunta e os trechos viram
vetores, e entram no contexto os trechos mais parecidos com a pergunta. Essa estratégia
funciona quando um único trecho contém a resposta, mas falha em perguntas
\textit{multi-hop}, que exigem encadear fatos espalhados por documentos diferentes. Na
pergunta ``Que prêmio recebeu quem criou a linguagem que influenciou o C++?'', o trecho
que fala do prêmio de Dennis Ritchie não menciona o C++ e, por isso, dificilmente está
entre os trechos mais parecidos com a pergunta.

As abordagens GraphRAG representam o conhecimento extraído dos documentos como um grafo
de entidades e relações e usam a estrutura desse grafo na recuperação
\cite{edge2024graphrag,peng2024graphragsurvey}. O KAG (\textit{Knowledge Augmented
Generation}) vai além e propõe a indexação mútua entre o grafo e os trechos de texto, o
raciocínio guiado por formas lógicas e o alinhamento do conhecimento extraído
\cite{liang2024kag}.

Este trabalho apresenta um sistema GraphRAG/KAG de perguntas e respostas sobre a
história da computação, construído a partir de artigos da Wikipédia em português. As
contribuições são: (i) um modelo de grafo de propriedades rotuladas no Neo4j com duas
camadas, texto e domínio, em que cada relação guarda o trecho que a sustenta; (ii) um
pipeline Text-to-Graph que extrai entidades e relações com um LLM restrito a um schema
de domínio; (iii) uma recuperação híbrida que combina busca vetorial, expansão
multi-hop e consultas Cypher geradas por LLM, com citação dos caminhos e trechos e um
nível de confiança calculado a partir de sinais observáveis; (iv) uma interface que
mostra o subgrafo e a procedência de cada resposta; e (v) uma avaliação comparativa entre
o RAG vetorial e o GraphRAG.

O restante do artigo está organizado assim. A Seção~\ref{sec:relacionados} discute
trabalhos relacionados. A Seção~\ref{sec:arquitetura} apresenta a arquitetura. As
Seções~\ref{sec:modelagem} a \ref{sec:interface} detalham a modelagem do grafo, o
pipeline, a recuperação e a interface. A Seção~\ref{sec:avaliacao} apresenta a
avaliação, e a Seção~\ref{sec:conclusao} conclui.

% =====================================================================
\section{Trabalhos relacionados}
\label{sec:relacionados}

\citeasnoun{lewis2020rag} propuseram combinar um recuperador denso com um gerador, de
modo que a resposta dependa de documentos recuperados e não apenas dos parâmetros do
modelo. O levantamento de \citeasnoun{ji2023survey} organiza os tipos de alucinação e
aponta o apoio em fontes externas como uma das principais formas de mitigação.

A combinação de LLMs com grafos de conhecimento é discutida por \citeasnoun{pan2024unifying},
que descrevem tanto o uso de grafos para dar fatos verificáveis aos LLMs quanto o uso de
LLMs para construir grafos. O GraphRAG de \citeasnoun{edge2024graphrag} extrai um grafo de
entidades dos documentos, agrupa as entidades em comunidades com o algoritmo de Leiden e
gera resumos dessas comunidades para responder a perguntas globais sobre o corpus. O
levantamento de \citeasnoun{peng2024graphragsurvey} classifica as variantes de GraphRAG
pelas etapas de indexação, recuperação e geração.

O KAG \cite{liang2024kag} foi proposto para domínios profissionais, em que a precisão
importa mais do que a cobertura. Suas ideias centrais são a indexação mútua entre o grafo
e os trechos de origem, o uso de formas lógicas para decompor a pergunta em passos de
consulta e o alinhamento semântico do conhecimento extraído.

Este trabalho adota a indexação mútua do KAG, ligando cada entidade aos trechos que a
mencionam e cada relação ao trecho que a sustenta, e usa medidas de grafo como a
centralidade do PageRank \cite{page1999pagerank} e as comunidades do Louvain
\cite{blondel2008louvain}. O foco é a resposta local a perguntas multi-hop, com
procedência até o texto original e recusa explícita quando o contexto não tem a resposta.

% =====================================================================
\section{Arquitetura}
\label{sec:arquitetura}

A Figura~\ref{fig:arquitetura} mostra os módulos do sistema. O pipeline transforma os
artigos em arquivos JSON com trechos, entidades e relações. O ingestor grava esses
arquivos no Neo4j. A recuperação responde às perguntas consultando o grafo e um LLM. A
API expõe a recuperação por HTTP, e a interface mostra a resposta e as evidências.

\begin{figure}[ht]
\centering
\begin{tikzpicture}[
  caixa/.style={draw, rounded corners, align=center, font=\footnotesize,
                minimum height=0.9cm, text width=2.2cm, fill=gray!8},
  seta/.style={-{Stealth[length=2mm]}, thick}]
  \node[caixa] (wiki) {Wikipédia\\(32 artigos)};
  \node[caixa, right=0.5cm of wiki] (chunk) {Chunking\\(trechos)};
  \node[caixa, right=0.5cm of chunk] (extr) {Extração\\(LLM + schema)};
  \node[caixa, right=0.5cm of extr] (ing) {Ingestor\\(validação)};
  \node[caixa, below=1.1cm of ing, fill=blue!8] (neo) {Neo4j\\(grafo + índices)};
  \node[caixa, left=0.5cm of neo] (rec) {Recuperação\\(vetor + grafo)};
  \node[caixa, left=0.5cm of rec] (api) {API\\(FastAPI)};
  \node[caixa, left=0.5cm of api] (ui) {Interface\\(Streamlit)};
  \draw[seta] (wiki) -- (chunk);
  \draw[seta] (chunk) -- (extr);
  \draw[seta] (extr) -- (ing);
  \draw[seta] (ing) -- (neo);
  \draw[seta] (neo) -- (rec);
  \draw[seta] (rec) -- (api);
  \draw[seta] (api) -- (ui);
\end{tikzpicture}
\caption{Módulos do sistema, da coleta dos artigos à interface.}
\label{fig:arquitetura}
\end{figure}

Dois contratos ligam os módulos. O primeiro é o schema do domínio, um arquivo JSON com os
rótulos, as relações e suas direções, a regra dos identificadores e o formato do arquivo
por artigo. O mesmo arquivo alimenta o prompt de extração e o prompt de geração de
Cypher, o que evita que os dois módulos usem nomes diferentes para a mesma relação. O
segundo é o contrato da resposta, definido com Pydantic, com o texto, a confiança, a
consulta Cypher, os caminhos, o subgrafo, os nós consultados, as fontes e os avisos.

O banco roda em contêiner Docker, com o Neo4j 5.26 (LTS, edição Community) e os plugins
APOC e Graph Data Science (GDS) embutidos numa imagem própria, com versões fixas e o
download do plugin conferido por \textit{checksum}. Um único comando sobe o banco e
aplica constraints e índices.

% =====================================================================
\section{Modelagem do grafo}
\label{sec:modelagem}

O conhecimento é representado como um grafo de propriedades rotuladas (LPG), em que nós
e relações têm rótulos e propriedades \cite{robinson2015graphdatabases}. O grafo tem duas
camadas ligadas entre si (Figura~\ref{fig:schema}).

\begin{figure}[ht]
\centering
\begin{tikzpicture}[
  no/.style={draw, rounded corners, align=center, font=\scriptsize,
             minimum height=0.8cm, text width=2.3cm},
  lex/.style={no, fill=gray!10},
  dom/.style={no, fill=blue!8},
  rel/.style={-{Stealth[length=1.8mm]}, thick},
  rot/.style={font=\scriptsize, fill=white, inner sep=1pt}]
  \node[lex] (doc) {:Document\\id, title, source};
  \node[lex, right=2.4cm of doc] (c1) {:Chunk\\id, index, text,\\embedding};
  \node[lex, right=2.4cm of c1] (c2) {:Chunk};
  \node[dom, below=1.5cm of doc] (pes) {:Entity:Pessoa\\Dennis Ritchie};
  \node[dom, right=2.4cm of pes] (tec) {:Entity:Tecnologia\\C};
  \node[dom, right=2.4cm of tec] (org) {:Entity:Organizacao\\Bell Labs};
  \node[dom, below=1.3cm of pes] (pre) {:Entity:Premio\\Prêmio Turing};
  \node[dom, below=1.3cm of tec] (cpp) {:Entity:Tecnologia\\C++};
  \draw[rel] (c1) -- node[rot, above] {PART\_OF} (doc);
  \draw[rel] (c1) -- node[rot, above] {NEXT} (c2);
  \draw[rel] (c1) -- node[rot, pos=0.4] {MENTIONS} (pes);
  \draw[rel] (c1) -- node[rot, pos=0.5] {MENTIONS} (tec);
  \draw[rel] (pes) -- node[rot, above] {CRIOU} (tec);
  \draw[rel] (tec) -- node[rot, above, font=\tiny] {DESENVOLVIDA\_EM} (org);
  \draw[rel] (pes) -- node[rot, left] {RECEBEU} (pre);
  \draw[rel] (cpp) -- node[rot, right] {INFLUENCIADA\_POR} (tec);
\end{tikzpicture}
\caption{Camada léxica (cinza) e camada de domínio (azul). Toda relação de domínio
guarda em \texttt{source\_chunk\_ids} os trechos que a sustentam.}
\label{fig:schema}
\end{figure}

A \textbf{camada léxica} contém os documentos (\texttt{Document}) e seus trechos
(\texttt{Chunk}), cada um com o texto e o vetor de embedding. A relação \texttt{PART\_OF}
liga o trecho ao documento, e \texttt{NEXT} liga trechos consecutivos. A \textbf{camada de
domínio} contém as entidades extraídas. Toda entidade tem o rótulo genérico
\texttt{Entity} e um rótulo do domínio. As camadas se ligam de duas formas: a relação
\texttt{MENTIONS} vai do trecho às entidades que ele cita, e cada relação de domínio
guarda na propriedade \texttt{source\_chunk\_ids} os identificadores dos trechos que a
sustentam. Essa indexação mútua, inspirada no KAG \cite{liang2024kag}, permite ir de um
trecho às suas entidades e de uma relação ao texto que a justifica, o que é a base da
procedência das respostas.

\subsection{Schema do domínio}

O domínio usa quatro rótulos e cinco relações com direção fixa
(Tabela~\ref{tab:dominio}). Optou-se por um schema pré-definido em vez de deixar o LLM
criar tipos livremente: isso reduz a variação entre artigos e torna previsíveis as
consultas Cypher geradas na recuperação.

\begin{table}[ht]
\centering
\caption{Rótulos e relações do domínio.}
\label{tab:dominio}
\footnotesize
\begin{tabular}{lll}
\toprule
\textbf{Relação} & \textbf{De $\rightarrow$ Para} & \textbf{Exemplo} \\
\midrule
\texttt{CRIOU} & Pessoa $\rightarrow$ Tecnologia ou Organizacao & Dennis Ritchie criou o C \\
\texttt{TRABALHOU\_EM} & Pessoa $\rightarrow$ Organizacao & Ritchie trabalhou no Bell Labs \\
\texttt{DESENVOLVIDA\_EM} & Tecnologia $\rightarrow$ Organizacao & Unix desenvolvido no Bell Labs \\
\texttt{INFLUENCIADA\_POR} & Tecnologia $\rightarrow$ Tecnologia & C++ influenciado pelo C \\
\texttt{RECEBEU} & Pessoa $\rightarrow$ Premio & Ritchie recebeu o Prêmio Turing \\
\bottomrule
\end{tabular}
\end{table}

Cada entidade tem um identificador canônico no formato \texttt{prefixo:nome}, em que o
nome é normalizado: sem a desambiguação entre parênteses, com ``+'' trocado por ``mais'',
sem acentos, em minúsculas e com os demais caracteres trocados por sublinhado. Assim,
``C++'' vira \texttt{tecnologia:c\_mais\_mais}. É esse identificador que faz a mesma
entidade, citada em artigos diferentes, virar um único nó.

\subsection{Restrições e índices}

A Tabela~\ref{tab:indices} resume as restrições e os índices, conforme listados por
\texttt{SHOW CONSTRAINTS} e \texttt{SHOW INDEXES} no banco carregado (todos com estado
\texttt{ONLINE}). As restrições de unicidade sobre o \texttt{id} impedem nós duplicados.
Os índices Range atendem a filtros e ordenações; no Neo4j 5, o índice BTree pedido no
enunciado foi substituído pelo Range. Os índices fulltext usam o analisador
\texttt{brazilian}, adequado ao português, e os vetoriais usam 384 dimensões e
similaridade de cosseno, com busca aproximada por HNSW \cite{malkov2020hnsw}.

\begin{table}[ht]
\centering
\caption{Restrições e índices do banco.}
\label{tab:indices}
\footnotesize
\begin{tabular}{l>{\raggedright\arraybackslash}p{5.4cm}>{\raggedright\arraybackslash}p{4.2cm}}
\toprule
\textbf{Tipo} & \textbf{Nome} & \textbf{Alvo} \\
\midrule
Unicidade & \texttt{document\_id\_unique}, \texttt{chunk\_id\_unique}, \texttt{entity\_id\_unique} & \texttt{id} de Document, Chunk e Entity \\
Range & \texttt{entity\_name}, \texttt{entity\_type} & \texttt{name} e \texttt{type} de Entity \\
Range & \texttt{entity\_pagerank}, \texttt{entity\_community} & resultados do GDS \\
Range & \texttt{document\_source} & URL do documento \\
Fulltext & \texttt{entity\_search} & \texttt{name}, \texttt{aliases}, \texttt{description} \\
Fulltext & \texttt{chunk\_text} & texto do trecho \\
Vetorial & \texttt{chunk\_embedding} & embedding do trecho (384, cosseno) \\
Vetorial & \texttt{entity\_embedding} & embedding da entidade (384, cosseno) \\
\bottomrule
\end{tabular}
\end{table}

A edição Community do Neo4j não oferece restrições de existência de propriedade nem
usuários com papéis restritos. Por isso, a obrigatoriedade das propriedades é verificada
pelo ingestor, e a proteção contra escrita nas consultas geradas pelo LLM é feita
executando-as em transação de leitura, em que o banco recusa qualquer escrita.

\subsection{Ingestão}

O ingestor lê um arquivo JSON por artigo e valida cada um contra o schema do domínio:
chaves obrigatórias, tipos de entidade e de relação, direção das relações, prefixo dos
identificadores e existência dos trechos referenciados. Em seguida, calcula o embedding de
cada trecho com o mesmo módulo usado depois para as perguntas e grava o artigo numa
única transação: ou o artigo entra inteiro, ou nada dele entra. Toda escrita usa
\texttt{MERGE} pelo identificador, de modo que reprocessar a mesma pasta não altera as
contagens.

\subsection{Algoritmos de grafo e o grafo carregado}

Depois da ingestão, o GDS projeta em memória apenas as entidades e as relações entre elas
e executa o PageRank \cite{page1999pagerank} e o Louvain \cite{blondel2008louvain}. O
PageRank, que convergiu dentro do limite de 100 iterações, é gravado em cada entidade e usado pela
recuperação para desempatar caminhos; as comunidades do Louvain são gravadas para
enriquecer a análise e a visualização.

A Tabela~\ref{tab:grafo} resume o grafo carregado com os 17 artigos já processados pelo
pipeline até a escrita deste texto, e a Tabela~\ref{tab:pagerank} lista as entidades de
maior PageRank. Das 265 comunidades encontradas pelo Louvain, 27 têm mais de uma entidade,
e as maiores correspondem a temas reconhecíveis: Tim Berners-Lee, o CERN e a Web; a família
C, C++, Java e Simula; a Sun Microsystems; e Dennis Ritchie, Ken Thompson, o Bell Labs e o
Unix. Nenhuma relação ficou sem trecho de origem. O relatório de validação também apontou
pontos a melhorar na extração: 9 pares de prováveis duplicatas (por exemplo, ``MIT'' e
``Instituto de Tecnologia de Massachusetts''), 238 das 467 entidades sem nenhuma relação
de domínio e 228 dos 567 trechos sem nenhuma entidade mencionada.

\begin{table}[ht]
\centering
\caption{Grafo carregado com os 17 artigos processados.}
\label{tab:grafo}
\small
\begin{tabular}{ccccc}
\toprule
\textbf{Documentos} & \textbf{Trechos} & \textbf{Entidades} & \textbf{Relações} & \textbf{Comunidades} \\
\midrule
17 & 567 & 467 & 259 & 265 \\
\bottomrule
\end{tabular}
\end{table}

\begin{table}[ht]
\centering
\caption{As cinco entidades de maior PageRank.}
\label{tab:pagerank}
\small
\begin{tabular}{clll}
\toprule
\textbf{\#} & \textbf{Entidade} & \textbf{Rótulo} & \textbf{PageRank} \\
\midrule
1 & Tim Berners-Lee & Pessoa & 11,24 \\
2 & Sun Microsystems & Organizacao & 10,51 \\
3 & Linus Torvalds & Pessoa & 5,99 \\
4 & Alan Kay & Pessoa & 5,95 \\
5 & James Gosling & Pessoa & 5,59 \\
\bottomrule
\end{tabular}
\end{table}

% =====================================================================
\section{Pipeline Text-to-Graph}
\label{sec:pipeline}

\textbf{Coleta.} O dataset parte de uma lista de 32 artigos da Wikipédia em português,
escolhidos por se citarem entre si: pioneiros e o Prêmio Turing, a família C e Unix,
Java e Sun, Python, Lisp e linguagens históricas, Linux e a web. Um script baixa o texto
puro de cada artigo, sem marcação e sem as seções de referências, e registra num
manifesto o título, a URL, a revisão exata, a data de acesso e a licença (CC BY-SA 4.0).

\textbf{Divisão em trechos.} O modelo de embeddings lê no máximo 128 \textit{tokens} de
cada texto. Por isso, os artigos são divididos em trechos de até 450 caracteres, cortados
entre frases, com a última frase de cada trecho repetida no início do seguinte. Só uma
frase que sozinha passa do limite é quebrada, entre palavras. Com essa regra, 17 dos
1.633 trechos do dataset (1\%) ainda passam de 128 \textit{tokens}.

\textbf{Extração.} As entidades e relações são extraídas por um LLM, o modelo
\texttt{openai/gpt-oss-120b}, acessado pela API do Groq. Os trechos de cada artigo são
enviados em lotes de até cinco, e o prompt traz o schema do domínio. A resposta do modelo é
validada contra o domínio, e o que não respeita os tipos ou as direções é descartado. As
respostas ficam em cache, para que uma nova execução não repita chamadas já feitas, e os
limites de uso da API são tratados com novas tentativas espaçadas.

\textbf{Resolução de entidades.} A resolução atribui o identificador canônico e une
entidades equivalentes a partir de mapeamentos explícitos. Uma triagem semântica
conservadora sinaliza relações suspeitas para revisão manual, sem removê-las nem
invertê-las automaticamente. Até a escrita deste texto, 17 dos 32 artigos tinham passado
pela extração e pela resolução, e todos os 17 arquivos passaram na validação do ingestor.

% =====================================================================
\section{Recuperação}
\label{sec:recuperacao}

A recuperação recebe uma pergunta e funciona em dois modos. No modo vetorial, usa apenas
os trechos mais parecidos com a pergunta. No modo GraphRAG, parte desses trechos, chega às
entidades que eles citam e percorre o grafo a partir delas. Nos dois modos, a resposta é
gerada por um LLM a partir de um contexto numerado, as citações são conferidas, a
confiança é calculada e o resultado é validado contra o contrato da resposta. Todas as
consultas ao Neo4j rodam em transação de leitura.

\textbf{Busca vetorial.} Trechos e perguntas são convertidos em vetores normalizados de 384
dimensões pelo modelo \texttt{paraphrase-multilingual-MiniLM-L12-v2}
\cite{reimers2019sbert,reimers2020multilingual}. A busca devolve os cinco trechos mais
parecidos, com o título e a URL do documento. O índice devolve o score na faixa de 0 a 1,
calculado como $(1+\cos)/2$, e o cálculo da confiança desfaz essa conversão. Com o grafo
semente, um conjunto de trechos, entidades e relações anotados à mão, 12 de 12 perguntas
de teste tiveram um trecho correto entre os três primeiros resultados.

\textbf{Expansão multi-hop.} A partir das entidades citadas nos trechos recuperados, a
expansão percorre de 1 a 3 saltos, usando só as relações do domínio. Cada caminho é escrito
como uma frase e comparado com a pergunta pelo mesmo modelo de embeddings. O sistema
identifica pelo início da pergunta o tipo de entidade pedido (``Quem'' pede uma pessoa,
``Que prêmio'' pede um prêmio) e ordena os caminhos dando prioridade aos que terminam
nesse tipo; depois, pela similaridade com a pergunta, pelo score do trecho de origem, pela
média do PageRank dos nós e pelo tamanho. Os cinco primeiros caminhos entram no contexto,
junto com os trechos que sustentam suas relações, até o limite de oito fontes.

\textbf{Text-to-Cypher.} No modo GraphRAG, o sistema também pede ao LLM uma consulta Cypher.
O prompt traz o schema do domínio, os identificadores das entidades encontradas na pergunta
pelo índice fulltext e três exemplos. A consulta gerada recebe \texttt{LIMIT 25} se não
tiver limite, passa pelo \texttt{EXPLAIN} do Neo4j, que detecta erros sem executá-la, e
só então roda, em transação de leitura. Se falhar, o LLM recebe a mensagem de erro e tem
mais uma tentativa. A proteção contra escrita fica com o banco: em transação de leitura,
o Neo4j recusa \texttt{CREATE}, \texttt{MERGE}, \texttt{SET} e \texttt{DELETE}. Os caminhos da
Cypher disputam lugar com os da expansão na mesma ordenação, e uma falha do Text-to-Cypher
vira um aviso, sem interromper a resposta.

\textbf{Grounding e procedência.} O LLM recebe um contexto em que cada caminho é
identificado como C1, C2\ldots{} e cada trecho como F1, F2\ldots, e deve responder só com
base nele, citando cada afirmação. Se o contexto não tiver a resposta, deve escrever
``Não encontrei essa informação no dataset.'' A chamada usa a API da Abacus RouteLLM com
temperatura 0. Citações inexistentes não são apagadas: viram avisos na resposta. Cada
caminho guarda os trechos que sustentam suas relações, e cada fonte traz o trecho, o
documento e a URL, de modo que uma afirmação pode ser seguida até o artigo de origem.

\textbf{Nível de confiança.} A confiança sai só de sinais observáveis da recuperação, sem
pedir ao LLM que avalie a própria resposta. No modo GraphRAG, é a média entre o cosseno
do melhor trecho e a maior similaridade entre a pergunta e um caminho, multiplicada pela
fração de caminhos com trecho de origem em todas as relações; uma falha do Text-to-Cypher
desconta 0,1. No modo vetorial, é 0,7 vezes o cosseno, pois sem caminhos não há relações a
conferir. O valor é zero quando não há contexto ou quando o modelo recusa. A partir de
0,7 o nível é alto, de 0,4 a 0,7 é médio e, abaixo disso, baixo.

\textbf{Exemplo.} Para a pergunta ``Que prêmio recebeu quem criou a linguagem que
influenciou o C++?'', o Text-to-Cypher gerou e executou a consulta:

{\small
\begin{verbatim}
MATCH p = (:Tecnologia {id: 'tecnologia:c_mais_mais'})
    -[:INFLUENCIADA_POR]->(:Tecnologia)<-[:CRIOU]-(:Pessoa)
    -[:RECEBEU]->(:Premio)
RETURN p LIMIT 25
\end{verbatim}
}

A resposta citou dois caminhos, como o LLM os recebeu no contexto:

{\small
\begin{verbatim}
[C1] C++ -[INFLUENCIADA_POR]-> C <-[CRIOU]- Dennis Ritchie
     -[RECEBEU]-> Prêmio Turing
[C4] C++ -[INFLUENCIADA_POR]-> C <-[CRIOU]- Dennis Ritchie
     -[RECEBEU]-> Medalha Nacional de Tecnologia e Inovação
\end{verbatim}
}

O texto gerado afirma que o C++ foi influenciado pelo C, criado por Dennis Ritchie no
AT\&T Bell Labs, e que Ritchie recebeu o Prêmio Turing em 1983. O trecho que sustenta o
prêmio não estava entre os cinco da busca vetorial: entrou nas fontes porque sustenta a
relação \texttt{RECEBEU} do caminho C1. No modo vetorial, nenhum dos cinco trechos do
contexto falaria de prêmio. A confiança foi 0,59 (nível médio): cosseno do melhor trecho
0,64, similaridade do caminho 0,54 e cinco de cinco caminhos com trecho de origem.

% =====================================================================
\section{Interface e explicabilidade}
\label{sec:interface}

A interface foi construída com Streamlit e consome a API, que expõe a rota
\texttt{POST /perguntar} com a pergunta e o modo (\texttt{graphrag} ou \texttt{vetorial}) e
devolve o objeto do contrato da resposta. Além do texto da resposta, a interface mostra:

\begin{itemize}
\item o \textbf{subgrafo de contexto}, desenhado com a biblioteca streamlit-agraph
(baseada em vis-network), com nós arrastáveis, cores por tipo de entidade, tamanho
proporcional ao PageRank, rótulos nas relações, destaque dos caminhos citados e detalhes do
nó selecionado;
\item a \textbf{lista de nós consultados};
\item a \textbf{consulta Cypher} gerada e executada;
\item o \textbf{nível de confiança}, acompanhado do critério com os números usados no
cálculo;
\item as \textbf{fontes}, com o trecho, o documento e a URL, e os \textbf{avisos}, como
citações inexistentes ou falhas do Text-to-Cypher.
\end{itemize}

Com esses painéis, o usuário pode conferir de onde veio cada afirmação, em vez de confiar
apenas no texto gerado.

% =====================================================================
\section{Avaliação}
\label{sec:avaliacao}

A avaliação comparou o RAG vetorial e o GraphRAG em 15 perguntas com gabarito: 5 de um
salto, 7 multi-hop e 3 sem resposta no dataset. Cada pergunta tem a resposta esperada, os
trechos que a sustentam e, nas multi-hop, o caminho esperado no grafo. Foram medidas: a
\textit{evidência recuperada}, quando todos os trechos do gabarito estão nas fontes; o
\textit{caminho completo}, quando todas as relações do gabarito estão no subgrafo; e a
\textit{recusa correta}, nas perguntas sem resposta. Cada resposta foi ainda classificada
manualmente como correta, incompleta ou alucinada. A Tabela~\ref{tab:avaliacao} resume os
resultados.

\begin{table}[ht]
\centering
\caption{RAG vetorial e GraphRAG nas 15 perguntas de avaliação.}
\label{tab:avaliacao}
\small
\begin{tabular}{llcccccc}
\toprule
\textbf{Modo} & \textbf{Tipo} & \textbf{N} & \textbf{Evid.} & \textbf{Caminho} & \textbf{Correta} & \textbf{Incompl.} & \textbf{Alucin.} \\
\midrule
Vetorial & Um salto & 5 & 0\% & 0\% & 0 & 5 & 0 \\
Vetorial & Multi-hop & 7 & 0\% & 0\% & 1 & 6 & 0 \\
Vetorial & Sem resposta & 3 & -- & -- & 3 & 0 & 0 \\
GraphRAG & Um salto & 5 & 0\% & 0\% & 0 & 5 & 0 \\
GraphRAG & Multi-hop & 7 & 14\% & 29\% & 2 & 5 & 0 \\
GraphRAG & Sem resposta & 3 & -- & -- & 3 & 0 & 0 \\
\bottomrule
\end{tabular}
\end{table}

Nas perguntas multi-hop, o GraphRAG deu duas respostas corretas em sete, e o vetorial,
uma. A diferença está na pergunta sobre o prêmio de quem criou a linguagem que deu origem
ao C: o modo vetorial recusou, porque a busca não trouxe o trecho do prêmio, que não cita
o C nem o B, enquanto o GraphRAG chegou a esse trecho seguindo o grafo, do C ao B, do B a
Ken Thompson e dele ao Prêmio Turing. O GraphRAG também montou o caminho completo do
gabarito em duas das sete perguntas multi-hop. Os dois modos recusaram as três perguntas
sem resposta, e nenhuma resposta foi classificada como alucinada: quando o contexto não
tinha a resposta, o sistema recusou em vez de inventar.

Os resultados têm um limite claro. Quando as respostas foram geradas, o banco tinha
apenas quatro documentos, e 20 das 24 respostas a perguntas com resposta são recusas de
perguntas que dependem de artigos que não estavam carregados. Por isso, nenhum dos modos
acertou as perguntas de um salto, e os números indicam uma tendência, não uma diferença
estatisticamente significativa. A análise dos erros também mostrou casos em que a
expansão parou antes do último salto, por falta de uma relação extraída, e um caso em que
a consulta Cypher filtrou uma entidade pelo nome em vez do identificador e não encontrou
caminho.

% =====================================================================
\section{Conclusão}
\label{sec:conclusao}

Este trabalho apresentou um sistema GraphRAG/KAG de perguntas e respostas sobre a
história da computação. O grafo no Neo4j combina uma camada léxica, com os trechos e seus
embeddings, e uma camada de domínio, com entidades e relações restritas a um schema
pré-definido, ligadas por indexação mútua. Essa ligação permitiu que a recuperação chegasse
a trechos que a busca vetorial não encontraria e que cada afirmação da resposta fosse
seguida até o texto original. Na avaliação, o GraphRAG respondeu corretamente uma pergunta
multi-hop a mais que o RAG vetorial, um indício que precisa ser confirmado com o corpus
completo, e nenhum dos modos produziu resposta alucinada.

As principais limitações são o dataset ainda parcial, com 17 dos 32 artigos processados
pela extração, as duplicatas e lacunas apontadas pelo relatório de validação, a checagem
de citações, que confere se o identificador existe, mas não se o trecho sustenta a
afirmação, e a dependência do LLM na geração das consultas Cypher. Como trabalhos futuros,
ficam o processamento de todo o dataset, uma resolução de entidades mais ampla, o uso de
embeddings das entidades na recuperação, a verificação de que cada citação sustenta a
afirmação e a repetição da avaliação com os 17 artigos já carregados e com mais perguntas.

% =====================================================================
\section*{Créditos}

O código está no repositório do projeto, com instruções de reprodução. O sistema usa o
Neo4j 5.26 Community com os plugins APOC e Graph Data Science; o driver Python do Neo4j;
a biblioteca sentence-transformers e o modelo
\texttt{paraphrase-multilingual-MiniLM-L12-v2}, sobre o PyTorch; a API do Groq, com o
modelo \texttt{openai/gpt-oss-120b}, na extração; a API da Abacus RouteLLM na geração das
respostas e das consultas Cypher; FastAPI, Uvicorn e Pydantic na API; Streamlit e
streamlit-agraph na interface; pytest nos testes; e Docker no ambiente do banco.

Os textos do dataset são da Wikipédia em português e estão sob a licença Creative Commons
Atribuição-CompartilhaIgual 4.0 (CC BY-SA 4.0). Os autores são os editores de cada artigo,
listados no histórico das páginas. O título, a URL e a revisão exata de cada artigo usado
estão no manifesto do dataset, no repositório.

\IfFileExists{sbc.bst}{\bibliographystyle{sbc}}{\bibliographystyle{apalike}}
\bibliography{referencias}

\end{document}
